\documentclass[11pt]{scrartcl}
\usepackage[utf8]{inputenc}
\usepackage[sexy]{evanchen} 
\usepackage{graphicx}
\graphicspath{{8 Outsource K Math/Gambar/}}

\title{Power of Point}
\author{Azzam and Akeyla}

\begin{document}

\maketitle

\section{Definisi}
Power of Point (POP) dari sebuah titik $X$ terhadap lingkaran $\omega$ yang berpusat di $O$ dan berjari-jari $r$ adalah \[\text{pow}_\omega(X) = OX^2 - r^2\]

\section{Pengaruh Posisi Titik Terhadap Lingkaran}
\subsection{Titik $X$ berada di luar lingkaran}
Ikuti langkah-langkah berikut.
\begin{enumerate}
    \item Tarik garis singgung dari titik $X$ ke lingkaran $\omega$ yang menyinggung di titik $P$ dan $Q$. Jelas bahwa $XP = XQ$. Buktikan bahwa $\text{pow}_\omega(X) = XP^2 = XQ^2$.
    \item Tarik garis yang dari $X$ yang memotong $\omega$ di titik $A$ dan $B$. Buktikan bahwa segitiga $XAP$ sebangun dengan $XPB$.
    \item Buktikan juga bahwa $\text{pow}_\omega(X) = XA\cdot XB$
    \item Simpulkan bahwa apabila ditarik dua garis dari titik $X$ yang memotong $\omega$ masing-masing di titik $A, B$ dan $C, D$, maka $XA\cdot XB = XC\cdot XD$. Pernyataan ini benar bahkan jika $C = D$, yakni saat $XC$ adalah garis singgung $\omega$
    \item Buktikan inversnya: Apabila terdapat segiempat konveks $ABCD$ dengan $AB$ dan $CD$ berpotongan di $X$ dan $XA\cdot XB = XC\cdot XD$, maka $ABCD$ merupakan segiempat siklik. 
\end{enumerate}
\subsection{Titik $X$ berada di keliling lingkaran}
Yakinkan diri kalian bahwa $\text{pow}_{\omega}(X) = 0$

\subsection{Titik $X$ berada di dalam lingkaran}
Ikuti langkah-langkah berikut.
\begin{enumerate}
    \item Periksa apakah power dari $X$ bernilai positif, negatif, atau nol.
    \item Tarik garis $OX$ yang memotong lingkaran $\omega$ di dua titik $A$ dan $B$
    \item Melambangkan apakah nilai $AX \cdot XB$?
    \item Tarik tali busur lain yang melewati titik $X$. Misalkan tali busur ini memotong di titik $C$ dan $D$. Buktikan bahwa segitiga $AXC$ dan $BXD$ sebangun.
    \item Buktikan bahwa $\text{pow}_{\omega}(X) = - XC\cdot XD$
    \item Simpulkan bahwa apabila ditarik dua garis dari titik $X$ yang memotong $\omega$ masing-masing di titik $A, B$ dan $C, D$, maka $XA\cdot XB = XC\cdot XD$.
    \item Buktikan inversnya: Apabila terdapat segiempat konveks $ABCD$ dengan $AC$ dan $BD$ berpotongan di $X$ dan $XA\cdot XC = XB\cdot XD$, maka $ABCD$ merupakan segiempat siklik. 
\end{enumerate}

\section{Radical Axis}
Misalkan terdapat dua buah lingkaran nonkonsentris (alasannya?) $\omega_1$ dan $\omega_2$. Radical axis dari kedua lingkaran ini didefinisikan sebagai lokus semua titik yang memiliki power yang sama ke kedua lingkaran ini.

\begin{enumerate}
    \item (Opsional) Buktikan bahwa lokus tersebut merupakan sebuah garis yang tegak lurus terhadap garis yang menghubungkan kedua titik pusat lingkaran tersebut.
    \item Tentukan radical axis dari dua buah lingkaran apabila kedua lingkaran tersebut berpotongan sebanyak: 
    \begin{enumerate}
        \item $2$ kali
        \item $1$ kali
        \item (Opsional) $0$ kali
    \end{enumerate}
    \item \textbf{Radical Center.} Jika terdapat tiga lingkaran dengan titik pusat yang tidak segaris, buktikan bahwa tiga radical axis yang terbentuk akan berpotongan di satu titik. Titik perpotongan ini disebut sebagai \textit{radical center}. Bagaimana jika ketiga titik pusat lingkaran tersebut segaris?
    \item \textbf{Degenerate Circle.} Power dan Radical axis juga dapat didefinisikan apabila lingkaran yang diacu merupakan sebuah titik. Hal ini dikarenakan titik dapat dianggap sebagai lingkaran berjari-jari nol. Bagaimana bentuk radical axis dari:
    \begin{enumerate}
        \item Lingkaran dan titik?
        \item Titik dan titik?
    \end{enumerate}
\end{enumerate}

\section{Latihan}
\subsection{Cek Ombak}
\begin{problem}
    Diberikan sebuah lingkaran dengan tali busur $AB$ dan $CD$ yang berpotongan di $X$. Jika diketahui bahwa $AX = 2, BX = 4, CX = 1$, panjang dari $DX$ adalah ...
\end{problem}

\subsection{Mudah}
\begin{problem}
    (AMC 2020 12B P10) Di sebuah persegi satuan $ABCD$, lingkaran dalam $\omega$ menyinggung $CD$ di titik $M$. Jika $AM$ memotong $\omega$ di titik $P$, panjang dari $AP$ adalah ...
\end{problem}

\begin{problem}
    Misalkan $\omega_1$ dan $\omega_2$ adalah dua lingkaran yang berpotongan di dua titik $A$ dan $B$. Garis $\ell$ adalah garis singgung persekutuan kedua lingkaran. $\ell$ menyinggung lingkaran  masing-masing di titik $P$ dan $Q$. Buktikan bahwa perpanjangan garis $AB$ membagi dua segmen garis $PQ$ sama panjang.
\end{problem}

\subsection{Sedang}
\begin{problem}
    Diberikan segitiga $ABC$. Lingkaran dalam segitiga $ABC$ menyinggung sisi $BC, CA, AB$ berturut-turut di titik $D, E, F$. Titik $X$ dan $Y$ berturut-turut adalah titik tengah $BD$ dan $DC$. Titik $P$ dan $Q$ berturut-turut adalah titik tengah $BF$ dan $CE$. Buktikan bahwa $XP$ dan $YQ$ berpotongan di garis sumbu sisi $BC$.
\end{problem}

\begin{problem}
    Diberikan dua buah lingkaran $\omega_1$ dan $\omega_2$ yang berpotongan di titik $S$ dan $T$. Garis singgung persekutuan kedua lingkaran menyinggung kedua lingkaran masing-masing di $A$ dan $B$. Jika $AS = x, BS = y,$ dan $AT = z$, tentukan panjang $BT$.
\end{problem}

\begin{problem}
    (Swedia MO 2019 P2) $AB$ adalah diameter dari sebuah lingkaran. Titik $C$ dan $D$ terletak di keliling lingkaran. Sinar $AC$ dan $AD$ secara berturut-turut memotong garis singgung lingkaran dari titik $B$ di titik $P$ dan $Q$. Buktikan bahwa $P, C, D, Q$ membentuk segiempat siklik.     
\end{problem}

\begin{problem}
    Diberikan setengah lingkaran yang berdiameter $AB$. Titik $C$ terletak di keliling setengah lingkaran (bukan di diameter). Titik $E$ adalah proyeksi titik $D$ terhadap garis $BC$. Titik $F$ adalah perpotongan $AE$ dengan setengah lingkaran. Buktikan bahwa $BF$ memotong segmen garis $DE$ sama panjang.
\end{problem}

\begin{problem}
    (Irlandia MO 2019 P5) Titik $M$ terletak di sisi $BC$ dari sebuah segitiga $ABC$. Titik $P$ dan $Q$ secara berturut-turut adalah titik pusat lingkaran luar $ABM$ dan $ACM$. Titik $L$ merupakan perpotongan $BP$ dan $CQ$. Titik $K$ merupakan pencerminan titik $L$ terhadap garis $PQ$. Buktikan bahwa titik $M, P, Q,$ dan $K$ terletak di keliling lingkaran yang sama.
\end{problem}

\begin{problem}
    (USAMO 1998) Misal lingkaran $\omega_1$ dan $\omega_2$ adalah lingkaran yang konsentris dengan $\omega_2$ terletak di interior $\omega_1$. Titik $A$ terletak di keliling $\omega_1$ dan $B$ terletak di keliling $\omega_2$ sedemikian hingga $AB$ merupakan garis singgung $\omega_2$. $C$ adalah titik perpotongan kedua $AB$ $\omega_1$. Titik $D$ adalah titik tengah $AB$. Sebuah garis yang melewati $A$ dan memotong $\omega_2$ di titik $E$ dan $F$ sedemikian hingga garis sumbu $DE$ dan garis sumbu $CF$ berpotongan di titik $M$ yang terletak di segmen garis $AB$. Tentukan rasio dari $AM:MC$.
\end{problem}

\begin{problem}
    Diberikan segitiga $ABC$ dengan titik pusat lingkaran dalam $I$, titik pusat lingkaran luar $O$, jari-jari lingkaran dalam $r$, dan jari-jari lingkaran luar $R$. Buktikan bahwa $OI^2 = R(R - 2r)$
\end{problem}

\end{document}